\documentclass[../../../include/open-logic-section]{subfiles}

\begin{document}

\olfileid{sth}{replacement}{strength}
\olsection{د تعويض قوت}

د تعويض د قوت په اړه له يوې ساده کتنې پيل کوو: که له $\Z$ څخه
وړاندې ولاړ نۀ شو، د $\omega + \omega$ سره برابر يا تر هغه لوی هېڅ د
فون نويمن ترتيبي عدد شتون نۀ شو ثابتولے۔

د دې دليل لنډه طرحه داسې ده۔ په~$\ZF$ کښې کار وکړئ او
$V_{\omega+\omega}$ سټ په پام کښې ونيسئ۔ دا سټ د~$\Z$ د يوه
\emph{مدل} د ساحې کار کوي۔ د دې د ليدلو دپاره د فارمول د
\emph{نسبي کولو} نښه تعريفوو:
\begin{defn}\ollabel{formularelativization}
	د هر سټ $M$ او هر فارمول $\phi$ دپاره، $\phi^M$ هغه فارمول دے
چې د $\phi$ ټول کميت ټاکونکي په $M$ پورې محدود کړو۔ يعنې
``$\lexists[x]$'' په ``$(\lexists[x \in M])$'' او
``$\lforall[x]$'' په ``$(\lforall[x \in M])$'' بدل کړئ۔
\end{defn}\noindent
ښودل کېدے شي چې د هر اصل $\phi$ د~$\Z$ دپاره
$\ZF \vdash \phi^{V_{\omega+\omega}}$ دے۔ خو د
\olref[spine][rank]{ordsetrankalpha} له مخې $\omega+\omega$ په
$V_{\omega+\omega}$ کښې \emph{غړے نۀ دے}۔ نو $\Z$ د
$\omega+\omega$ له نشتون سره سازګار دے۔

له همدې امله مو په \olref[ordinals][replacement]{sec} کښې ويلي وو
چې \olref[ordinals][ordtype]{thmOrdinalRepresentation} له تعويض پرته
نۀ شي ثابتېدلے۔ ځکه په~$\Z$ کښې يو څرګند ښه ترتيب په اسانۍ تعريفولے
شو چې په بداهتي توګه ئې د ترتيب ډول بايد $\omega+\omega$ وي۔ د دې
يوه غير رسمي مثال مو په \olref[ordinals][idea]{sec} کښې د طبيعي
عددونو په دې ترتيب سره ورکړے و:
\begin{align*}
	n \lessdot m \text{ هله او يوازې هله چې }&\text{يا }n < m\text{ او }m-n\text{ جفت وي،}\\
	& \text{يا $n$ جفت او $m$ طاق وي.}
\end{align*}
خو که $\omega+\omega$ موجود نۀ وي، دا ښه ترتيب له هېڅ ترتيبي عدد
سره آيزومورف نۀ دے۔ نو $\Z$ د
\olref[ordinals][ordtype]{thmOrdinalRepresentation} ثبوت \emph{نۀ}
ورکوي۔

اوس موضوع بل اړخ ته واړوو: تعويض موږ ته د $\omega+\omega$ شتون
ثابتوي؛ نو د $V_{\omega+\omega}$ شتون هم ورسره ثابتولے شي۔ دا يوازې
همدومره نۀ دي۔ د هر هغه ښه ترتيب دپاره چې تعريفولے ئې شو،
\olref[ordinals][ordtype]{thmOrdinalRepresentation} وايي چې له هغه
ښه ترتيب سره آيزومورف يو $\alpha$ شته، او له دې امله $V_\alpha$ هم
شته۔ نو تعويض په څرګنده توګه دا تضمينوي چې د سټونو سلسله بايد
\emph{ډېره لوړه} وي۔

په راتلونکو څو برخو کښې، او بيا په
\olref[card-arithmetic][fix]{sec} کښې، به لا ښه پوه شو چې تعويض د
سټونو سلسله \emph{څومره} لوړېدو ته اړ باسي۔ اوس دپاره ساده خبره دا ده
چې تعويض په رښتيا \emph{توجيه} غواړي!

\end{document}
